\documentclass[11pt,a4paper]{article}
\usepackage{xeCJK}
\usepackage{fontspec}
\usepackage{xcolor}
\usepackage{fancyhdr}
\usepackage{booktabs}
\usepackage{array}
\usepackage{amssymb}
\usepackage{geometry}
\usepackage[protrusion=true]{microtype}
\geometry{margin=1.8cm,top=2cm,headheight=15pt}
\linespread{1.05}

\setmainfont{Baskerville}
\setCJKmainfont{Songti SC}
\newfontfamily{\starfont}{Songti SC}

\definecolor{wrongred}{RGB}{198,40,40}
\definecolor{wrongbg}{RGB}{253,235,233}
\definecolor{rulegray}{RGB}{200,200,200}

\setlength{\emergencystretch}{3em}
\setlength{\fboxsep}{3pt}
\setlength{\parindent}{0pt}

\newcommand{\checkmarkbox}{\fbox{\rule{0pt}{0.8ex}\rule{0.8ex}{0pt}}}
\newcommand{\STAR}{\textcolor{wrongred}{\starfont ★}}
\newcommand{\bk}{\rule[-0.55ex]{1.0em}{0.4pt}}
\newcommand{\A}[1]{\textbf{#1}}
\newcommand{\W}[1]{\textcolor{wrongred}{#1}}
\newcommand{\OK}{\textcolor{wrongred}{\ensuremath{\checkmark}}}
\newcommand{\NO}{\textcolor{wrongred}{$\times$}}
\newcolumntype{L}[1]{>{\raggedright\arraybackslash}p{#1}}

\pagestyle{fancy}
\fancyhf{}
\fancyhead[L]{\footnotesize 英语语法 时态语态练习1（填空专项 · 订正）}
\fancyhead[R]{\footnotesize 赵文瀚}
\fancyfoot[C]{\thepage}
\renewcommand{\headrulewidth}{0.4pt}
\renewcommand{\footrulewidth}{0pt}

\begin{document}

\begin{center}
{\LARGE\bfseries 时态语态练习1 \quad 填空专项}\\[0.25em]
{\large 错题订正与考点复盘}\\[0.35em]
{\footnotesize 2029 届高一英语 · 智学网 AI 错题本（原错题 1--18）· 2026-09}
\end{center}

\vspace{-0.2em}
\noindent{\color{rulegray}\rule{\textwidth}{0.8pt}}

\vspace{0.5em}
\noindent\fcolorbox{black!15}{black!4}{%
  \parbox{\dimexpr\linewidth-2\fboxsep-2\fboxrule\relax}{%
    \small 日期：2026-09\qquad 来源：智学网 AI 错题本「时态语态练习1」（共 18 空）\qquad 姓名：赵文瀚\\[0.3em]
    条目：\textbf{18}\,/\enspace 18（原错题本收录，均为原卷失分点）\qquad
    重做日期：\underline{\hspace{2.4cm}}}}

\vspace{0.6em}
\noindent{\footnotesize 说明：来源（智学网 AI 错题本）只给了「原错题 + 参考答案」，\textbf{未给学生原始作答}，故本表不设「我的作答」栏，改为「正确答案 + 考点／易错点」。第 9 题答案册写 \emph{were working}，按标准语法 \emph{everybody} 作单数，本表取 \emph{was working}（见脚注）。}

\vspace{0.4em}
\renewcommand{\arraystretch}{1.06}
\noindent{\footnotesize\begin{tabular}{@{}L{0.72cm}@{\hspace{0.08cm}}L{9.2cm}@{\hspace{0.08cm}}L{3.1cm}@{\hspace{0.08cm}}L{3.6cm}@{}}
\toprule
\textbf{题号} & \textbf{题干（含空与原提示词）} & \textbf{正确答案} & \textbf{考点 / 易错点} \\
\midrule
1  & This time next month I \bk\ (sit) on the shuttle bus. & \A{will be sitting} & 将来进行（this time next month） \\
2  & Yesterday was the second time his brother \bk\ (be) late for school. & \A{had been} & 过去完成（It was the second time…） \\
3  & Once he \bk\ (make) up his mind to do something, you will never hold him back. & \A{has made} & 现在完成表将来（once…） \\
4  & The students \bk\ (write) busily when Miss Brown went to get a book that she \bk\ (leave) in the office. & \A{were writing; had left} & 过去进行 + 过去完成（一句两空） \\
5  & My mother \bk\ (work) as a teacher for 30 years by the end of this year. & \A{will have worked} & 将来完成（by the end of…） \\
6  & Betty, \bk\ (always make) the same mistake. I need to talk to her. & \A{is always making} & 现在进行表不满（always…） \\
7  & The recent years, \bk\ (see) great changes in Shanghai. & \A{have seen} & 现在完成（In recent years） \\
8  & Every means \bk\ (try) since then. & \A{has been tried} & 现在完成的被动（means 单数） \\
9  & Everybody in the office \bk\ (work) hard when the manager came in. & \A{was working} & 过去进行；everybody 单数 \\
10 & I don't think Jim saw me; he \bk\ (stare) into space. & \A{was staring} & 过去进行（当时正在） \\
11 & ---Hi, Tracy, you look tired. ---I am tired. I \bk\ (paint) the living room all day. & \A{have been painting} & 现在完成进行（all day） \\
12 & Shirley \bk\ (write) a book about China last year but I don't know whether she has finished it. & \A{was writing} & 过去进行（去年那段时间） \\
13 & ---Oh, it's you! I \bk\ (not recognize) you. ---I \bk\ (just have) my hair cut, and I \bk\ (wear) new glasses. & \A{didn't recognize; have just had; am wearing} & 一句三空：一般过去 + 现在完成 + 现在进行 \\
14 & We \bk\ (publish) ten books in the last three years. & \A{have published} & 现在完成（in the last three years） \\
15 & It's said that the scientist \bk\ (work) in Africa with wild animals for six years in the 1990's. & \A{worked} & 一般过去（in the 1990's） \\
16 & I don't really work in this shop. I \bk\ (help) out here these days. & \A{am helping} & 现在进行表临时（these days） \\
17 & He said in the email that he \bk\ (leave) for the international conference the next day. & \A{was leaving / would leave} & 过去将来（the next day） \\
18 & I \bk\ (think) of going up to Scotland at the end of December, but I've only got a few days' holiday. & \A{was thinking} & 过去进行（原打算） \\
\bottomrule
\end{tabular}}

\vspace{0.3em}
\noindent{\footnotesize\textbf{原卷旁注}：（一）第 9 题答案册作 \emph{were working}；按标准语法 \emph{everybody} 作主语是\textbf{单数}，应为 \emph{was working}，请以教材口径为准。（二）第 11 题原卷题号前残留「31 ……」，系源练习编号。（三）原卷有两处笔误：\emph{The stcdents}（应 students）、\emph{with wild animal}（应 animals）。（四）第 13 题第一个空用一般过去 \emph{didn't recognize}（说话瞬间才认出），第二个空用 \emph{have just had}（刚刚做完），第三个空用 \emph{am wearing}（佩戴状态）。}

\newpage
\begin{center}
{\Large\bfseries 考点归类与复盘}
\end{center}
\vspace{-0.4em}
\noindent{\color{rulegray}\rule{\textwidth}{0.8pt}}

\vspace{0.5em}
\noindent 说明：按语法考点归类 18 个空（第 4、13 题为一句多空，单列）。

\vspace{0.5em}
\renewcommand{\arraystretch}{1.3}
\begin{center}
\begin{tabular}{@{}p{5.6cm} p{8.4cm} c@{}}
\toprule
\textbf{考点} & \textbf{题号} & \textbf{题数} \\
\midrule
现在完成时 / 现在完成进行 & {\small 3, 7, 8, 11, 14} & 5 \\
过去进行时 & {\small 9, 10, 12, 18} & 4 \\
将来时（将来进行、将来完成） & {\small 1, 5} & 2 \\
现在进行时（表临时 / 不满） & {\small 6, 16} & 2 \\
一句多空综合 & {\small 4, 13} & 2 \\
过去完成时 & {\small 2} & 1 \\
过去将来时 & {\small 17} & 1 \\
一般过去时 & {\small 15} & 1 \\
\midrule
\textbf{合计} & \textbf{18 空} & \textbf{18} \\
\bottomrule
\end{tabular}
\end{center}

\vspace{0.5em}
\noindent\textbf{关联知识点标签}
\vspace{0.25em}

\noindent\begin{tabular}{@{}p{8.2cm}@{\hspace{0.4cm}}p{8.2cm}@{}}
\begin{minipage}[t]{\linewidth}
\small
\STAR\ \textbf{时间状语定调}：this time next month（将来进行）、by the end of…（将来完成）、in the last/recent years、since then（现在完成）、in the 1990's（一般过去）\\[0.2em]
\STAR\ \textbf{"第二次"句型}：It/This is the first time (that) sb \textbf{has done}；It was the second time sb \textbf{had done}\\[0.2em]
\STAR\ \textbf{all day / 一直}：现在完成进行 have been doing
\end{minipage} &
\begin{minipage}[t]{\linewidth}
\small
\STAR\ \textbf{when 从句背景}：主句常用过去进行（were writing / was working）\\[0.2em]
\STAR\ \textbf{现在进行两用}：表"临时"（am helping out these days）／表"情绪"（is always making）\\[0.2em]
\STAR\ \textbf{过去将来}：主句过去 + the next day $\to$ was leaving / would leave\\[0.2em]
\STAR\ \textbf{被动与单数一致}：Every means has been tried；everybody + 单数谓语
\end{minipage} \\
\end{tabular}

\vspace{0.65em}
\noindent\fcolorbox{black!15}{black!4}{%
  \parbox{\dimexpr\linewidth-2\fboxsep-2\fboxrule\relax}{%
    \textbf{复盘记录}\\[0.35em]
    最薄弱的板块：\underline{\hspace{9.0cm}}\\[0.55em]
    主要错误原因：\checkmarkbox 完成时/进行时混淆\quad \checkmarkbox 时间状语没抓准\quad \checkmarkbox 主谓一致\quad \checkmarkbox 被动语态\\[0.55em]
    本次复习的改进措施：\underline{\hspace{10.1cm}}\\[0.55em]
    \hspace{3.2em}\underline{\hspace{10.1cm}}\\[0.55em]
    下次复习日期：\underline{\hspace{2.2cm}}\qquad
    当前掌握度：\checkmarkbox 仍薄弱\quad \checkmarkbox 基本掌握\quad \checkmarkbox 已掌握
  }}

\end{document}
